\chapter{Central simple algebras and the Brauer group}
\label{ch:rm03-central-simple}

Throughout this chapter \(k\) is a field. Algebras and algebra maps are
unital, and a simple algebra is nonzero and has no nonzero proper
two-sided ideal. For a field extension \(L/k\), write \(A_L=A\otimes_k L\).
We use the Artin--Wedderburn theorem from
\cite[ALG-03-C03-R04]{p103} and the finite-dimensional case of density
from Corollary~\ref{cor:rm03-density-finite}. For comparison with the
classical development, see \cite[Sections~11.3--11.8]{stackscsa}.
The finite and real division-ring inputs are the established results in
\cite[K07.23.finite\_division\_ring and K07.23.frobenius]{d07}.

\section{Structure and extension of scalars}

\begin{definition}[RM03-C04.01. Central simple algebras]
\label{def:rm03-csa}
A \emph{central simple \(k\)-algebra} is a finite-dimensional simple
\(k\)-algebra \(A\) whose center is \(k1_A\). A \emph{central division
algebra} over \(k\) is a finite-dimensional division algebra with this
same center condition. A central simple \(k\)-algebra is \emph{split}
if it is isomorphic to \(M_n(k)\) for some \(n\geq 1\). An extension
\(L/k\) \emph{splits} \(A\) if \(A_L\) is split over \(L\).
\end{definition}

\begin{proposition}[RM03-C04.02. Matrix and division presentations]
\label{prop:rm03-csa-matrix}
Every central simple \(k\)-algebra \(A\) admits a presentation
\(A\simeq M_r(D)\), where \(r\geq 1\) and \(D\) is a central division
algebra over \(k\). Both \(r\) and the \(k\)-algebra isomorphism class
of \(D\) are uniquely determined by \(A\).
\end{proposition}
\begin{proof}
A finite-dimensional simple algebra is left Artinian, so the simple
case of Artin--Wedderburn gives \(A\simeq M_r(D)\). Commuting with
the matrix units forces any central matrix to have the form \(dI_r\),
and commuting with coefficient matrices forces \(d\in Z(D)\).
Thus \(Z(D)=k\), with the identification induced by the original
\(k\)-action on \(A\). For a simple left \(A\)-module \(S\) the matrix
presentation gives \(S\simeq D^r\), with endomorphism ring \(D^{\mathrm{op}}\)
via right multiplication by \(D\). This identifies \(D\) as a
\(k\)-algebra since every module endomorphism commutes with the
given scalars. Finally \(r^2\dim_k D=\dim_k A\) determines \(r\).
\end{proof}

The same simple-Artinian structure will be used for a finite-dimensional
simple \(k\)-algebra \(T\) with arbitrary center. Write \(T=M_s(E)\)
for its Artin--Wedderburn presentation. A finite nonzero left
\(T\)-module \(V\) is a direct sum of \(m\geq 1\) copies of the
unique simple column module \(E^s\). In particular the
\(k\)-dimension of \(V\) determines this multiplicity. Its
endomorphism ring is \(M_m(E^{\mathrm{op}})\), by writing an
endomorphism in blocks between the \(m\) copies of \(E^s\).
Consequently \(\operatorname{End}_T(V)\) is simple, and
\begin{equation}
\label{eq:rm03-simple-end-dimension}
\dim_k T\;\dim_k\operatorname{End}_T(V)=(\dim_k V)^2.
\end{equation}
Indeed the two factors on the left are \(s^2\dim_k E\) and
\(m^2\dim_k E\), and \(\dim_k V=sm\dim_k E\).

\begin{proposition}[RM03-C04.03. The enveloping algebra of a central simple algebra]
\label{prop:rm03-enveloping}
If \(A\) is central simple over \(k\), its left and right multiplication
map is an isomorphism of \(k\)-algebras
\[
 A\otimes_k A^{\mathrm{op}}\longrightarrow \operatorname{End}_k(A),
 \qquad a\otimes b^{\mathrm{op}}\longmapsto (x\mapsto axb).
\]
\end{proposition}
\begin{proof}
The displayed map is an algebra homomorphism because composition
reverses the order of the right multipliers. Under this action the
left \(A\otimes A^{\mathrm{op}}\)-submodules of \(A\) are exactly its
two-sided ideals, so \(A\) is a simple module. Its endomorphisms as
such a module are determined by the image \(c\) of \(1\). Compatibility
with left and right multiplication gives \(ac=ca\) for every \(a\),
and each such \(c\) defines an endomorphism. The resulting endomorphism
division algebra is \(k\). Corollary~\ref{cor:rm03-density-finite} therefore
gives surjectivity onto \(\operatorname{End}_k(A)\). Both \(k\)-dimensions
are \((\dim_k A)^2\), which gives injectivity.
\end{proof}

\begin{proposition}[RM03-C04.04. Tensor product and base change]
\label{prop:rm03-csa-tensor}
Let \(A\) be central simple over \(k\) and \(B\) any simple
\(k\)-algebra. Then \(A\otimes_k B\) is simple, with center
\(1\otimes Z(B)\). In particular the tensor product of two central
simple \(k\)-algebras is central simple. Also \(A_L\) is central
simple over \(L\) for every field extension \(L/k\).
\end{proposition}
\begin{proof}
Let \(I\) be a nonzero two-sided ideal in \(A\otimes_k B\). Express
a nonzero \(u\in I\) as \(\sum_i a_i\otimes b_i\) in a basis
\((a_i)\) of \(A\), and choose \(j\) with \(b_j\ne 0\).
The ideal is preserved by left and right multiplication from
\(A\otimes 1\), hence is preserved by \(T\otimes\operatorname{id}_B\)
for every \(k\)-linear \(T:A\to A\) by
Proposition~\ref{prop:rm03-enveloping}. Take \(T(a_j)=1\) and
\(T(a_i)=0\) for \(i\ne j\). This gives \(1\otimes b_j\in I\).
Because \(Bb_jB\) generates the unit ideal of \(B\), left and
right multiplication by \(1\otimes B\) then gives \(1\otimes 1\in I\).

An element commuting with \(A\otimes 1\) can be written
\(\sum_j x_j\otimes y_j\) for \(k\)-linearly independent
\(y_j\in B\). Each \(x_j\) commutes with \(A\), so the element
belongs to \(1\otimes B\). An element of that subalgebra is central
exactly when it commutes with \(1\otimes B\), giving the center formula.
Applying it to \(B=L\), and noting
\(\dim_L A_L=\dim_k A\), gives the base change statement.
\end{proof}

\section{Separable splitting fields}

\begin{theorem}[RM03-C04.05. Degree and finite Galois splitting]
\label{thm:rm03-csa-galois-split}
For \(A\) central simple over \(k\), \(\dim_k A\) is a square.
Write \(\deg A=\sqrt{\dim_k A}\). There is a finite Galois extension
\(K/k\) with \(A_K\simeq M_{\deg A}(K)\).
More precisely, if \(A\simeq M_r(D)\) and \(d=\deg D\), then \(D\)
contains a separable field extension \(E/k\) of degree \(d\) which
splits \(D\). Thus \(\deg A=rd\).
\end{theorem}
\begin{proof}
Over an algebraic closure \(\overline{k}\), base change and
Proposition~\ref{prop:rm03-csa-matrix} give
\(A_{\overline{k}}\simeq M_n(D')\).
Every element \(x\) of \(D'\) satisfies a polynomial over
\(\overline{k}\) by finite dimensionality. Factoring the polynomial
into linear factors and using absence of zero divisors shows
\(x\in\overline{k}\). This proves \(A_{\overline{k}}\simeq
M_n(\overline{k})\) and \(\dim_k A=n^2\). The same argument for \(D\)
gives \(D_{\overline{k}}\simeq M_d(\overline{k})\) and
\(n=rd\).

It suffices to split \(D\), since matrix algebras commute with
extension of scalars. If \(k\) is finite, then \(D\) is a finite
division ring and is commutative by
\cite[K07.23.finite\_division\_ring]{d07}. Its center condition forces
\(D=k\), so \(E=k\) works.

Assume \(k\) is infinite. Fix a \(k\)-basis \(u_1,\ldots,u_{d^2}\)
of \(D\). Through the splitting over \(\overline{k}\), the matrix
corresponding to \(\sum_i X_i u_i\) has a characteristic polynomial
whose discriminant is a polynomial in the \(X_i\) over
\(\overline{k}\). This discriminant is nonzero as a polynomial:
the basis spans \(M_d(\overline{k})\) after extension, so choose
coefficients giving a diagonal matrix with \(d\) distinct entries.
Any nonzero polynomial over \(\overline{k}\) in finitely many
variables has a nonzero value at some tuple in \(k\). To see this,
view it as a polynomial in its last variable, choose the other
variables by induction so that at least one coefficient survives,
and then avoid the finitely many roots of the remaining polynomial.
Thus some \(x\in D\) corresponds to a matrix with \(d\) distinct
eigenvalues over \(\overline{k}\).

The minimal polynomial of \(x\) has the same degree after extension:
linear independence of the initial powers \(1,x,\ldots,x^{j-1}\)
is preserved, and a monic relation over \(k\) remains a relation
over \(\overline{k}\). Uniqueness of the monic relation of smallest
degree then identifies the two minimal polynomials. For the chosen
matrix this polynomial is the product of the \(d\) distinct linear
factors. Hence \(E=k[x]\) has dimension \(d\) and is separable over \(k\).
The algebra \(E\) is a field since it is a finite-dimensional
commutative domain inside \(D\).

Regard \(D\) as a right \(E\)-space of dimension \(d\). The map
\[
 D\otimes_k E\longrightarrow \operatorname{End}_E(D),\qquad
 a\otimes e\longmapsto (z\mapsto aze)
\]
is an \(E\)-algebra homomorphism: the right multipliers belong to
the commutative field \(E\). It is unital and its domain is simple
by Proposition~\ref{prop:rm03-csa-tensor}, so it is injective.
Both sides have \(E\)-dimension \(d^2\), giving a splitting over
\(E\). A finite Galois closure \(K\) of \(E/k\) then splits both \(D\)
and \(A\).
\end{proof}

\begin{corollary}[RM03-C04.06. Index and degrees of splitting extensions]
\label{cor:rm03-csa-index}
Let \(A\simeq M_r(D)\) be central simple and set
\(\operatorname{ind}(A)=\deg D\). Every finite splitting field
\(L/k\) of \(A\) has \(\operatorname{ind}(A)\mid [L:k]\).
A separable splitting field of degree \(\operatorname{ind}(A)\)
exists.
\end{corollary}
\begin{proof}
Write \(d=\deg D\). Since \(A_L=M_r(D_L)\) and \(D_L\) is central
simple by Proposition~\ref{prop:rm03-csa-tensor}, the uniqueness
of the division presentation implies that \(D_L\) is split
whenever \(A_L\) is split. Then \(D_L\simeq M_d(L)\), and
restricting its natural left action on \(L^d\) to \(D\) gives a
finite left \(D\)-space. Its \(k\)-dimension \(d[L:k]\) must be a
multiple of \(\dim_k D=d^2\). The last assertion is the field \(E\)
of Theorem~\ref{thm:rm03-csa-galois-split}.
\end{proof}

\section{Conjugacy and commutants}

\begin{theorem}[RM03-C04.07. Skolem--Noether]
\label{thm:rm03-skolem-noether}
Let \(A\) be central simple over \(k\), and let \(B\) be a
simple \(k\)-algebra. Two unital \(k\)-algebra embeddings
\(f,g:B\hookrightarrow A\) are conjugate by an element of \(A^\times\).
In particular every \(k\)-algebra automorphism of \(A\) is inner.
\end{theorem}
\begin{proof}
The embedding makes \(B\) finite-dimensional over \(k\), and
\(T=A\otimes_k B^{\mathrm{op}}\) is a finite-dimensional simple
\(k\)-algebra by Proposition~\ref{prop:rm03-csa-tensor}.
Give the \(k\)-space \(A\) two left \(T\)-module structures by
\[
 (a\otimes b^{\mathrm{op}})\cdot_f x=axf(b),\qquad
 (a\otimes b^{\mathrm{op}})\cdot_g x=axg(b).
\]
The use of \(B^{\mathrm{op}}\) compensates for the order of right
multiplication. The two modules have equal \(k\)-dimension, so the
simple-Artinian module description preceding
Proposition~\ref{prop:rm03-enveloping} gives a \(T\)-module
isomorphism \(F:(A,\cdot_f)\to(A,\cdot_g)\).
Because \(F\) is left \(A\)-linear, \(F(x)=xu\) for \(u=F(1)\).
Its inverse is also right multiplication, which shows that
\(u\in A^\times\). Intertwining the two actions of \(b^{\mathrm{op}}\)
on \(1\) gives \(f(b)u=ug(b)\). The case \(B=A\) and \(g\) the
identity gives the automorphism assertion.
\end{proof}

\begin{theorem}[RM03-C04.08. Centralizer and double centralizer]
\label{thm:rm03-centralizer}
Let \(B\) be a unital simple \(k\)-subalgebra of a central simple
\(k\)-algebra \(A\). Its centralizer
\[
 C_A(B)=\{\,x\in A:xb=bx\ \text{for every }b\in B\,\}=C
\]
is simple, \(\dim_k B\cdot\dim_k C=\dim_k A\), and \(C_A(C)=B\).
Moreover \(Z(C)=Z(B)\) as subalgebras of \(A\).
If \(B\) is central over \(k\), so is \(C\), and multiplication
induces an isomorphism \(B\otimes_k C\simeq A\).
\end{theorem}
\begin{proof}
Act on \(A\) by \(T=A\otimes B^{\mathrm{op}}\) using left and right
multiplication, where \(T\) is simple by
Proposition~\ref{prop:rm03-csa-tensor}. An endomorphism of this
\(T\)-module is left \(A\)-linear, hence is right multiplication
by some \(x\in A\). Right multiplication by \(x\) commutes with
right multiplication by every \(b\in B\) exactly when \(x\in C\).
Thus \(\operatorname{End}_T(A)\simeq C^{\mathrm{op}}\).
By \eqref{eq:rm03-simple-end-dimension}, \(C\) is simple and
\[
 (\dim_k A)(\dim_k B)(\dim_k C)=(\dim_k A)^2.
\]
The same dimension formula applied to the simple subalgebra
\(C\) gives \(\dim_k C_A(C)=\dim_k B\). The inclusion
\(B\subseteq C_A(C)\) therefore is an equality. Also
\(Z(C)=C\cap C_A(C)=C\cap B=Z(B)\).

If \(Z(B)=k\), multiplication is a unital homomorphism from
\(B\otimes_k C\) because the subalgebras commute. The tensor
product is simple by Proposition~\ref{prop:rm03-csa-tensor},
so the homomorphism is injective. The dimension formula gives
surjectivity.
\end{proof}

\begin{corollary}[RM03-C04.09. Maximal subfields of a central division algebra]
\label{cor:rm03-maximal-subfield}
Every maximal subfield \(E\) of a central division \(k\)-algebra
\(D\) has \([E:k]=\deg D\) and splits \(D\). The separable \(E\)
in Theorem~\ref{thm:rm03-csa-galois-split} is a maximal subfield.
\end{corollary}
\begin{proof}
For \(x\in C_D(E)\), the algebra \(E[x]\) is a commutative
finite-dimensional domain over \(E\), hence a field. Maximality
gives \(x\in E\), so \(C_D(E)=E\). Theorem~\ref{thm:rm03-centralizer}
now gives \([E:k]^2=\dim_k D\). The map using left multiplication
by \(D\) and right multiplication by \(E\) in the proof of
Theorem~\ref{thm:rm03-csa-galois-split} works for any \(E\) of
this degree. Conversely the same dimension formula forces
\(C_D(E)=E\) for every degree-\(\deg D\) subfield, which is
sufficient for maximality.
\end{proof}

\section{Reduced characteristic polynomial and norm}

\begin{proposition}[RM03-C04.10. Reduced characteristic polynomial]
\label{prop:rm03-reduced-norm}
Let \(A\) be central simple of degree \(n\) over \(k\).
For \(a\in A\), the characteristic polynomial of the image of
\(a\otimes 1\) in a matrix algebra over a splitting field belongs
to \(k[t]\) and is independent of the field and of the splitting
isomorphism. Denote it by \(\chi_a^{\mathrm{red}}(t)\) and set
\(\operatorname{Nrd}(a)=(-1)^n\chi_a^{\mathrm{red}}(0)\).
The reduced norm is a homogeneous polynomial of degree \(n\)
on \(A\), commutes with base change, is multiplicative, and
is nonzero exactly on \(A^\times\).
\end{proposition}
\begin{proof}
Choose a finite Galois splitting \(K/k\), an isomorphism
\(\phi:A_K\to M_n(K)\), and a \(k\)-basis \(u_1,\ldots,u_{n^2}\)
of \(A\). On the generic element put
\[
 P(t,\mathbf X)=
 \det\left(tI_n-\sum_i X_i\phi(u_i\otimes 1)\right).
\]
For \(\sigma\in\operatorname{Gal}(K/k)\), write \(\sigma_A\)
for the action on the coefficients of \(A_K\) and \(\sigma_M\)
for entrywise application to matrices. Then
\(\sigma_M\phi\sigma_A^{-1}\) is another \(K\)-algebra
isomorphism. By Theorem~\ref{thm:rm03-skolem-noether} the
two differ by a single conjugation on \(M_n(K)\). Hence
applying \(\sigma\) to the coefficients leaves \(P\) unchanged.
Consequently \(P\in k[t,\mathbf X]\). If \(L/k\) is another
splitting field, an algebraic closure of \(L\) contains a
\(k\)-copy of \(K\). Extending the two splitting isomorphisms
to this common field and using the same conjugacy argument
identifies their generic characteristic polynomials. These
identities also verify base change for a field extension
that does not itself split \(A\), by using a splitting field
over it and comparing there.

The constant term with sign \((-1)^n\) is the determinant
of the matrix of the generic element. Homogeneity and
multiplicativity follow from the determinant formulas after
extension to \(K\), and hence over \(k\). If \(a\) is invertible
then its matrix image is invertible. Conversely if its
matrix image is invertible, left multiplication \(x\mapsto ax\)
on \(A\) becomes invertible over \(K\). It is therefore
invertible over \(k\), giving some \(b\) with \(ab=1\).
Injectivity of left multiplication also gives \(ba=1\)
from \(a(ba-1)=0\).
\end{proof}

Over \(K\), the norm-one equation is the determinant-one
equation defining \(\operatorname{SL}_n\). Its descent is
\(\operatorname{SL}_1(A)\) in \cite[Result F05]{lag}, which
also gives the corresponding inner form \(\operatorname{PGL}_1(A)\).

\section{Brauer classes and examples}

\begin{theorem}[RM03-C04.11. Brauer equivalence and the group law]
\label{thm:rm03-brauer-group}
Two central simple \(k\)-algebras \(A,B\) are \emph{Brauer
equivalent} if \(M_p(A)\simeq M_q(B)\) for some \(p,q\geq 1\).
This is an equivalence relation, and each class has a unique
central division representative up to \(k\)-algebra isomorphism.
The classes form an abelian group \(\operatorname{Br}(k)\)
with laws
\[
 [A]+[B]=[A\otimes_k B],\qquad 0=[k],\qquad -[A]=[A^{\mathrm{op}}].
\]
The class \([A]\) is zero exactly when \(A\) is split. A field
extension \(L/k\) gives a homomorphism
\(\operatorname{Br}(k)\to\operatorname{Br}(L)\) by \([A]\mapsto[A_L]\).
\end{theorem}
\begin{proof}
Write \(A=M_r(D)\) and \(B=M_s(E)\) with division factors.
If they become isomorphic after passing to matrices, then
\(D\simeq E\) by Proposition~\ref{prop:rm03-csa-matrix}.
If \(D\simeq E\), both \(M_s(A)\) and \(M_r(B)\) are
isomorphic to \(M_{rs}(D)\). Thus equivalence means equality
of the division factors up to \(k\)-algebra isomorphism. The
class of \(k\) therefore consists exactly of split algebras.

The tensor product is central simple by
Proposition~\ref{prop:rm03-csa-tensor}. The identity
\(M_p(A)\otimes_k B\simeq M_p(A\otimes_k B)\), induced by
tensoring coefficients, makes the product independent of
the representative of \(A\). Commutativity via the tensor
flip handles the representative of \(B\). Associativity and
the identity class come from the tensor associator and unit.
The opposite algebra is central simple and
Proposition~\ref{prop:rm03-enveloping} identifies its product
with \(A\) as the split algebra \(\operatorname{End}_k(A)\).
Finally, extension of scalars commutes with the matrix and
tensor constructions and sends \(k\) to \(L\). These
identifications also give compatibility in towers of fields.
\end{proof}

\begin{example}[RM03-C04.12. The real and quaternion classes]
\label{ex:rm03-brauer-real}
There are exactly two real Brauer classes, represented by
\(\mathbb R\) and \(\mathbb H\). To see this, apply
Frobenius's real algebraic division-ring theorem
\cite[K07.23.frobenius]{d07} to the central division
representative. Finite-dimensionality guarantees
algebraicity, and \(\mathbb C\) is excluded because its
center is larger than \(\mathbb R\). The quaternion class
is nonzero by uniqueness of the division representative.
Quaternion conjugation \(q\mapsto\bar q\) reverses products,
so \(\mathbb H^{\mathrm{op}}\simeq\mathbb H\) and
\(2[\mathbb H]=0\). Thus
\(\operatorname{Br}(\mathbb R)\simeq\mathbb Z/2\mathbb Z\).

More explicitly \(\mathbb C=\mathbb R[i]\) is the degree-two
splitting subfield of \(\mathbb H\) from
Corollary~\ref{cor:rm03-maximal-subfield}, so
\(\mathbb H_{\mathbb C}\simeq M_2(\mathbb C)\).
Write \(q=a+bi+cj+dk\) in Hamilton's basis, with
\(i^2=j^2=-1\) and \(k=ij=-ji\). In a complex splitting one
can take
\[
 i\longmapsto \begin{pmatrix}\mathrm i&0\\0&-\mathrm i\end{pmatrix},
 \qquad j\longmapsto \begin{pmatrix}0&1\\-1&0\end{pmatrix},
 \quad \mathrm i^2=-1\ \text{in }\mathbb C .
\]
These matrices satisfy the quaternion relations, giving a
unital homomorphism out of \(\mathbb H_{\mathbb C}\).
It is an isomorphism by simplicity and equal dimension.
Its determinant on \(q\) is \(a^2+b^2+c^2+d^2\), recovering
the quaternion norm of \cite[K07.23.quaternion\_algebra]{d07}.
\end{example}

\begin{remark}[RM03-C04.13. Relative classes and descent]
\label{rem:rm03-relative-brauer}
For \(K/k\) finite Galois, let
\(\operatorname{Br}(K/k)=\ker(\operatorname{Br}(k)\to\operatorname{Br}(K))\).
It consists of the classes split by \(K\), by
Theorem~\ref{thm:rm03-brauer-group}. Put \(G=\operatorname{Gal}(K/k)\).
In fixed degree \(n\), the algebras split by \(K\) up to
isomorphism correspond to \(H^1(G,\operatorname{PGL}_n(K))\)
with its coefficientwise action. We first give the descent argument.

For a finite-dimensional \(K\)-space \(V\) with a semilinear
\(G\)-action \(T_\sigma\), the fixed vectors span \(V\) over \(K\).
Otherwise a nonzero \(K\)-linear functional \(\ell\) annihilates
their \(K\)-span. The vector
\(\sum_{\sigma\in G} T_\sigma(\lambda v)\) is fixed for each
\(v\in V\), \(\lambda\in K\). Applying \(\ell\) gives
\(\sum_\sigma\ell(T_\sigma v)\sigma(\lambda)=0\) for all \(\lambda\).
Distinct field automorphisms are linearly independent as
\(K\)-valued functions: in a nonzero relation with the fewest
terms, evaluate at \(a\lambda\) and subtract a multiple of the
relation at \(\lambda\), choosing \(a\) on which two of the
automorphisms differ. This eliminates one term without
eliminating the other. A one-term relation is impossible by
evaluation at \(1\). Thus \(\ell(v)=0\) for all \(v\), a
contradiction. A \(K\)-basis of fixed vectors is a \(k\)-basis of
\(V^G\), since the coefficients of a fixed vector in this basis
are fixed by \(G\). Consequently \(K\otimes_k V^G\to V\) is an
isomorphism. If \(V\) is a \(K\)-algebra and the action respects
multiplication, this is an algebra isomorphism. A \(G\)-equivariant
\(K\)-algebra map descends by restriction to the fixed subalgebras.

Given a splitting of an algebra of degree \(n\), transport its
coefficient action to \(M_n(K)\). Composing each transported
action with inverse entrywise \(\sigma\) gives a \(K\)-algebra
automorphism, inner by Theorem~\ref{thm:rm03-skolem-noether}.
Two implementing invertible matrices differ by a scalar, since
the matrices commuting with all the matrix units are the scalar
ones. Thus the semilinear action gives a projective cocycle.
Conversely a projective cocycle acts on \(M_n(K)\) by semilinear
conjugation; the scalar defects disappear in conjugation. Its
fixed algebra is central simple of degree \(n\). Its scalar
extension is the matrix algebra, so any nonzero proper ideal
would give a nonzero proper ideal there by extending a vector-space
basis. A central element extends to a scalar matrix, and its
scalar is fixed by \(G\), hence lies in \(k\). Changing the
splitting conjugates the action by a \(K\)-algebra isomorphism,
which gives just the projective coboundary relation.
Conversely that relation gives an equivariant conjugation, which
restricts to an isomorphism on fixed algebras. This proves the
fixed-degree classification. Brauer equivalence allows
enlargement by matrices.

For the descent convention where the transported semilinear
action is \(x\mapsto g_\sigma \sigma(x)g_\sigma^{-1}\),
choose lifts \(G_\sigma\in\operatorname{GL}_n(K)\) of the
projective matrices \(g_\sigma\). Then
\(c_{\sigma,\tau}=G_\sigma\sigma(G_\tau)G_{\sigma\tau}^{-1}\)
is scalar by the projective cocycle identity. Associating the
product of three lifted maps in two ways gives
\[
 c_{\sigma,\tau}c_{\sigma\tau,\nu}
 =\sigma(c_{\tau,\nu})c_{\sigma,\tau\nu}.
\]
Changing each \(G_\sigma\) by a scalar changes \(c\) by a
coboundary. Changing the splitting conjugates the projective
datum by a coboundary, which one can lift as
\(G_\sigma\mapsto H^{-1}G_\sigma\sigma(H)\); this leaves the
scalar \(c_{\sigma,\tau}\) unchanged. Enlargement by a
split matrix algebra replaces \(G_\sigma\) by \(I_m\otimes G_\sigma\)
and preserves \(c\). Thus the scalar class depends only
on the relative Brauer class. Tensoring lifts multiplies
the scalar cocycles. We have constructed the group homomorphism
\[
 \operatorname{Br}(K/k)\longrightarrow
 H^2(\operatorname{Gal}(K/k),K^\times),\qquad [A]\longmapsto[c].
\]
With the same lifts, the cocycle in
\cite[Equation~(59.61.5.1)]{stacksdescent} is \(c^{-1}\).
The homomorphism is injective. If the scalar class is zero,
rescale the \(G_\sigma\) so that \(c_{\sigma,\tau}=1\).
The operators \(v\mapsto G_\sigma\sigma(v)\) on \(K^n\) then
form a semilinear action. In a basis of invariant vectors the
action becomes entrywise \(\sigma\). The fixed endomorphisms
are the endomorphisms over \(k\) of the fixed vector space,
so the descended algebra is split.

For surjectivity take a normalized scalar \(2\)-cocycle \(c\).
On the \(K\)-space with basis \(e_\tau\), \(\tau\in G\), define
the invertible \(\sigma\)-semilinear operators
\[
 T_\sigma(e_\tau)=c_{\sigma,\tau}e_{\sigma\tau}.
\]
On \(e_\nu\), the cocycle identity gives
\(T_\sigma T_\tau=c_{\sigma,\tau}T_{\sigma\tau}\).
The operators therefore act genuinely on the full endomorphism
algebra by conjugation. The fixed algebra descends as above.
Relative to the displayed basis the \(T_\sigma\)'s are lifts
of its projective cocycle and the scalar defect is exactly \(c\).
This proves the relative isomorphism with the indicated sign.
Compare Colliot-Th\'el\`ene and Skorobogatov,
\href{https://www.imo.universite-paris-saclay.fr/~jean-louis.colliot-thelene/BGgroup_book.pdf}{Theorem~1.3.5}.

For the separable-closure statement, put \(G_k=\operatorname{Gal}(k^s/k)\)
and give \((k^s)^\times\) the discrete topology. A continuous cochain
on \(G_k^r\) has finite image and factors through a finite quotient:
choose product neighborhoods on which it is constant, cover the
compact domain by finitely many of these neighborhoods, and take a
common open normal subgroup refining all their factors. Enlarge the
finite Galois field fixed by that subgroup to contain the finitely
many cochain values. The cocycle equation, and any equality with a
coboundary, can then be checked on the finite quotient associated to
that enlarged field. This identifies continuous second cohomology
with the filtered colimit of the finite Galois groups'
second cohomology under inflation and coefficient inclusion.
The relative class maps agree with these transitions: extend the
splitting and use the same lifts inflated along restriction of
automorphisms. Every class of \(\operatorname{Br}(k)\) is split
at a finite Galois stage by Theorem~\ref{thm:rm03-csa-galois-split}.
Taking the union of relative classes therefore gives
\(\operatorname{Br}(k)\cong H^2_{\mathrm{cts}}(G_k,(k^s)^\times)\)
for this convention. Compare \cite[Theorem~59.61.6]{stacksdescent}.
\end{remark}
